\section{Introduction}
Scientific applications increasingly execute as dependency graphs of short
computations and intermediate objects. The performance problem extends beyond
placing a runnable task on a free core: a successor must find its inputs, the
runtime must preserve those inputs for other consumers or recovery, and a
worker must distinguish work that is waiting from work that can use a CPU.
As tasks become shorter, repeated coordination at these boundaries can occupy
a substantial part of the interval between useful computations.

In a conventional design, the same single-threaded manager handles both
logical scheduling (dependency checks, dispatch, retries, and completion) and
data-lifecycle coordination. A large workflow amplifies the latter demand:
each intermediate input or output can trigger identity lookups, peer-copy or
checkpoint admission, stage-in or stage-out coordination, acknowledgements,
and garbage-collection updates. The bytes may move directly between workers,
but the manager still serializes much of the control and lifetime bookkeeping.
When these events arrive faster than one reactor can drain them, the manager
becomes the bottleneck and workers can remain idle even when runnable work
exists.

Existing systems already exploit local intermediate storage and peer transfer.
TaskVine manages an ensemble's storage and transfers~\cite{taskvine}; its use in
high-energy physics demonstrates the importance of both data and function
execution costs~\cite{hep}. Consequently, a credible next step cannot be to
rediscover local caching or replace one copy path with another. The question is
which component should own each decision, and which events must precede a
successor's execution.

Consider a task whose producer has finished on worker A. Worker B may have an
execution opportunity before a backup copy has reached persistent storage.
Waiting for backup before making the successor eligible needlessly orders
recovery work before computation in the failure-free case. Yet immediately
preparing every eligible successor can fill B's disk or memory with work that
will wait in its queue. Increasing concurrent processes can hide some waits,
but a fixed factor can also intensify competition during compute-heavy phases.
These are related ordering and admission problems.

Our central design choice is to expose three boundaries: logical dispatch,
physical input preparation, and execution admission. DataVine's scheduler owns
DAG eligibility and task attempts; its data controller owns named intermediates,
replicas, and retention; workers own the final preparation and execution gates.
A task description can wait cheaply on a worker before its input sandbox is
prepared. A bounded queue absorbs dispatch latency while a separate execution
window responds to current CPU, runnable-process, and memory pressure.
The design preserves explicit coordination at ownership transitions rather than
claiming that independent components never communicate.

The research contribution is the concrete interaction of these boundaries for
file-producing scientific tasks. Consumer-driven data lookup, separated storage
planes, and local scheduling queues all have precedents~\cite{hoplite,pocket,apollo}.
DataVine does not claim to introduce those primitives. Its hypothesis is that
allowing a logically eligible task to be assigned before local input readiness,
while postponing preparation until execution approaches, can reduce coordination
waiting without staging the full dispatch queue. An aggregate execution window
then provides a second control over work that actually consumes resources.

This paper makes three contributions. First, it specifies an ownership and
event-ordering model that separates task success, replica availability, execution
admission, and requested-output durability. Second, it implements that model
in an existing runtime, including generation-bound queued-task recall, deferred
preparation, and opt-in attempt tracing. Third, it provides matched experiments
that hold task bodies and scientific checks constant while varying preparation,
window policy, and runtime. These include a complete real-data diphoton analysis,
expanded function-slot controls, and five-block experiments
across distinct physical hosts. The results distinguish gains of the implemented
execution path from unsupported claims that feedback or extra concurrency
alone explains them.

The central scalability consequence is that serialized logical scheduling does
not also serialize data-plane service work. The manager remains the single
owner of logical DAG state and dispatch, while Controller RPC handling and
durable-output persistence use independent concurrency pools. We measure the
resulting manager/data-plane separation explicitly and report the shared-state
boundary where metadata scaling stops.

The scope is intentionally testable. Current routing is local-first, then a
live peer, with an admitted controller backup as fallback. It is not a learned
three-way optimizer over peer, controller, and shared filesystem. Likewise,
the controller is a distinct owner but is not demonstrated to scale without
bound; worker overcommit is a heuristic with measurable failure modes. We
report these limits alongside the measurements so that the architecture can
be assessed independently of broader, unsupported novelty claims.
